\cleardoubleoddpage
\chapter{Discussion}
\label{chap:discussion}
\begingroup
\fontsize{11pt}{13.2pt}\selectfont

This chapter interprets the experimental results presented in
Chapter~\ref{chap:results}. We discuss the key findings for each track
separately, draw cross-domain insights from comparing the binary CIFAR-10
and inkjet quality control settings, and examine the practical and
theoretical implications of the results. Known limitations and directions
for future work conclude the chapter.

%% ==================================================================
\section{Key Findings: CIFAR-10 Binary CDM}
\label{sec:disc:cifar}

%% ------------------------------------------------------------------
\subsection{The Separation Loss Is the Dominant Performance Driver}
\label{sec:disc:cifar:separation}

The separation loss mainly acts as a stabiliser. The three-seed
$\lambda = 0$ baseline reaches $92.52\% \pm 11.07\%$ AUROC, but that
mean hides extreme seed sensitivity: the individual runs span
79.73\%, 98.81\%, and 99.01\%. Adding the separation loss moves the
mean to $99.03\% \pm 0.07\%$ at $\lambda = 0.02$, a smaller
+6.5\,pp gain than the seed-42-only comparison suggested, but with a
dramatic reduction in variance. The effect remains non-linear:
$\lambda = 0.001$ already yields 97.32\%, performance peaks at
$\lambda = 0.02$, and degrades at $\lambda = 0.1$ (96.67\%) as the
separation objective dominates the denoising loss.

%% ------------------------------------------------------------------
\subsection{Contrastive Scoring Is Essential}
\label{sec:disc:cifar:contrastive}

The OOD score must contrast reconstruction errors under both binary
conditions. Using only the ID-class error collapses to 20.2\% AUROC on
SVHN (essentially random) because global shifts in reconstruction
scale confound the absolute magnitude. The contrastive score
$s(x) = e_0(x) - e_1(x)$ cancels these confounders, measuring only the
model's relative preference for one condition over the other. Thus, the
ablation supports the need for contrastive use of the two conditions; it
does not by itself isolate an architecture with conditioning removed.

%% ------------------------------------------------------------------
\subsection{External OOD Generalisation Reveals Learned ID Manifold}
\label{sec:disc:cifar:external}

Across five auditable external benchmarks, AUROC ranges from 90.50\% to
96.97\%, with each benchmark evaluated against the full CIFAR-10 test
reference pool. Although CIFAR-100 is a near-OOD dataset (sharing low-level
visual statistics with CIFAR-10), it achieves the highest AUROC
(96.97\%), higher than far-OOD datasets such as Textures (92.84\%) and
SVHN (90.50\%). As noted in
Section~\ref{sec:meth:diffusion:formal},
this counterintuitive result arises because the airplane class is
under-represented in CIFAR-100, creating a strong distributional gap
that the CDM exploits effectively. The result demonstrates that semantic
proximity to the source dataset does not deterministically predict
detection difficulty; what matters is proximity to the specific ID class
manifold (airplane). SVHN and Textures are harder despite being far-OOD
because their low-level statistics (digit strokes, repeating textures)
overlap with certain airplane image regions. Food-101 and STL-10 are
retained as legacy traceability values only.

%% ------------------------------------------------------------------
\subsection{Ablation Insights: Efficiency and Design Choices}
\label{sec:disc:cifar:ablation}

\paragraph{Monte Carlo trials.}
Performance improves rapidly to $K{=}10$ then saturates
(Section~\ref{sec:results:ablations}),
making $K{=}10$ the practical recommendation. The $K$-AUROC curve was
benchmarked at $\lambda{=}0.01$; transferability of this saturation
point to the $\lambda{=}0.02$ regime was not independently verified.
Even $K{=}1$ provides useful signal, suggesting that a single denoising
step is a non-trivial distributional indicator.

\paragraph{Timestep sampling.}
Uniform sampling across $[1, T]$ is optimal. The per-timestep analysis
(Figure~\ref{fig:results:per-timestep-error}) shows OOD-discriminative
information distributed across the full noise range; concentrating on
intermediate timesteps discards useful signal from the tails.

\paragraph{Scoring method.}
Difference and ratio scoring perform comparably on the within-CIFAR
split, with ratio generalising slightly better to external datasets. We
retain difference scoring as the default for its lower FPR95 and simpler
interpretation.

%% ==================================================================

\section{Key Findings: Inkjet Print Quality Control}
\label{sec:disc:inkjet}

%% ------------------------------------------------------------------
\subsection{Public YOLO+CDM Performance}
\label{sec:disc:inkjet:public-cdm}

The retained public inkjet evidence shows that the YOLO+CDM pipeline is
a viable quality-classification signal on FTI\_Zer0P, reaching
$0.8673 \pm 0.0230$ AUROC at $\lambda = 0$ under 5-fold image-level
cross-validation. This is an encouraging result for a crop-based
generative classifier, but it should be interpreted within the released
evidence boundary: the thesis no longer uses unreleased multi-method
baseline comparisons as support for ranking claims.

%% ------------------------------------------------------------------
\subsection{Feature Difficulty Is Heterogeneous}
\label{sec:disc:inkjet:feature-difficulty}

Per-feature AUROC varies substantially across the public CDM folds
(Table~\ref{tab:results:inkjet-per-feature},
Figure~\ref{fig:results:per-feature-auroc}). Distance and dot features
are consistently easier, while edge roughness features are more variable.
This suggests that the local crop signal is sufficient for some geometric
defects but weaker for subtle texture or boundary defects.

%% ------------------------------------------------------------------
\subsection{Separation Loss Does Not Transfer Automatically}
\label{sec:disc:inkjet:separation}

Unlike the CIFAR-10 track, non-zero separation-loss weights do not
produce a statistically significant gain on the public inkjet folds
(Table~\ref{tab:results:inkjet-lambda}). The negative result is useful:
the loss can sharpen class-conditional reconstruction boundaries when
data scale and class separation are favourable, but it is not a universal
recipe for small, fine-grained industrial datasets.

%% ------------------------------------------------------------------
\subsection{YOLO Detection as a Pipeline Dependency}
\label{sec:disc:inkjet:yolo}

The two-stage pipeline introduces a detection dependency that is not
independently evaluated: we do not compare CDM classifications between
YOLO-detected and ground-truth crops. The YOLO model achieves 95.04\%
mAP@0.5, suggesting detection errors are infrequent but not negligible.
An ablation isolating the detection bottleneck is recommended for future
work.

%% ------------------------------------------------------------------
\section{Cross-Domain Insights}
\label{sec:disc:cross}

%% ------------------------------------------------------------------
\subsection{The Separation Loss Has Boundary Conditions}
\label{sec:disc:cross:separation}

The most important cross-domain finding is that the separation loss is
highly effective on CIFAR-10 but not significant on inkjet data
(Table~\ref{tab:results:cross-domain}).
The technique appears to depend on sufficient data scale and a clear
visual gap between conditions; in the public inkjet setting, the
available 1{,}204 images and subtle GOOD/BAD differences make the
separation signal less reliable than in CIFAR-10. This negative result
bounds the applicability of the technique and prevents
over-generalisation.

%% ------------------------------------------------------------------
\subsection{Contrastive Scoring Is Universal}
\label{sec:disc:cross:contrastive}

Both tracks confirm that the OOD/quality score must compare
reconstruction errors under two conditions. Contrastive scoring
normalises away global reconstruction scale differences and focuses on
the model's relative preference for one condition over the other, which is the
key insight separating conditional diffusion detectors from absolute
reconstruction-error methods.

%% ------------------------------------------------------------------
\subsection{Public Evidence Boundaries Matter}
\label{sec:disc:cross:public-evidence}

Both tracks reinforce a practical lesson about evidence quality. The
CIFAR-10 track supports broad ablation claims because the relevant
weights, scripts, and raw-score artefacts are public. The inkjet track
supports narrower claims: the public YOLO+CDM pipeline, per-feature
behaviour, and separation-loss ablation. Claims about broader method
rankings or representation choices require equally complete public
artefacts before they can be treated as retained thesis evidence.

%% ==================================================================

\section{Implications}
\label{sec:disc:implications}

%% ------------------------------------------------------------------
\subsection{Theoretical Implications}
\label{sec:disc:implications:theoretical}

Our results support conditional diffusion models as generative
classifiers: scoring via reconstruction-error contrast can be interpreted
as a likelihood-ratio-style comparison between competing conditions. The
auditable results are consistent with a manifold-based interpretation,
although detection difficulty is not explained by a simple near-vs-far
OOD ordering. The separation loss explicitly shapes the learnt manifold
to maximise the inter-condition reconstruction gap, a useful principle in
the large-scale CIFAR-10 setting studied here: when the training
objective does not directly optimise the evaluation metric, a targeted
auxiliary loss can improve the scoring signal.

%% ------------------------------------------------------------------
\subsection{Practical Implications}
\label{sec:disc:implications:practical}

The binary CDM with separation loss (${\sim}68.8$M parameters) achieves
99\% AUROC at $K = 50$ and 98.2\% at $K = 10$ (the $K$-AUROC curve was
benchmarked at $\lambda{=}0.01$; saturation behaviour at $\lambda{=}0.02$
was not independently verified). It is recommended when
high accuracy is required, inference cost is acceptable ($K = 10$:
about 3.0\,min per 10K images on an NVIDIA GeForce RTX 2080 Ti;
$K = 50$: about 21.6\,min), and the task has sufficient labelled data
with a clear visual separation signal. Simpler discriminative
alternatives are preferred for ultra-low latency or limited-data
settings. For
industrial QC specifically, the public inkjet results suggest that
additional data scale, feature-specific analysis, and careful validation
matter more than simply porting the CIFAR-10 separation-loss recipe.
For quality inspection pipelines, the practical recommendation is to
validate the public crop-based generative signal against the exact
deployment constraints before drawing broader architecture or
representation conclusions.

%% ==================================================================
\section{Limitations}
\label{sec:disc:limitations}

%% ------------------------------------------------------------------
\subsection{Experimental Limitations}
\label{sec:disc:limitations:experimental}

The evaluation is scoped to a single ID class (airplane), a single
benchmark dataset (CIFAR-10), a single industrial application (inkjet
QC), and external OOD datasets limited to natural images. Results should
therefore be interpreted as strong evidence within these two settings
rather than as universal claims. Extension to multi-class splits,
higher-resolution datasets, diverse industrial domains, and
adversarial/corrupted inputs would strengthen generalisability.

%% ------------------------------------------------------------------
\subsection{Methodological Limitations}
\label{sec:disc:limitations:methodological}

\paragraph{Inference cost.}
Even at $K{=}10$, the binary CDM is ${\sim}32{\times}$ slower than a
discriminative ResNet, precluding real-time applications.

\paragraph{Training cost.}
Beyond inference, the separation-loss objective itself imposes a
training-time overhead that is decoupled from the inference cost above.
Each optimisation step requires three forward passes of the UNet: one
for $\mathcal{L}_{\text{diffusion}}$ using the true label, plus two
additional passes under the fixed $c{=}0$ and $c{=}1$ conditions to form
$\mathcal{L}_{\text{sep}}$
(Equation~\ref{eq:meth:lsep}), so
training the binary CDM with $\lambda > 0$ costs approximately
$3{\times}$ the per-step compute of the $\lambda = 0$ baseline and is
correspondingly more expensive than a discriminative ResNet-18
fine-tuning run (which completes in ${\sim}2$--4 hours on an A100).
Practitioners evaluating the total computational investment of our method
should account for both the training overhead and the inference overhead:
the separation loss is cheap to add in code but not cheap to run, and
this cost is separate from the $K$-dependent inference budget.

\paragraph{Binary framing.}
The binary setup (one ID vs.\ one proxy class) simplifies the general
multi-class OOD problem; the proxy class is itself in-distribution data.

\paragraph{Separation loss tuning.}
No principled method exists for selecting $\lambda$; it must be treated
as a task-dependent hyper-parameter selected via ablation.

\paragraph{Threshold calibration.}
All methods require threshold calibration assuming access to
representative OOD or defective samples during validation.

%% ------------------------------------------------------------------
\subsection{Theoretical Limitations}
\label{sec:disc:limitations:theoretical}

\paragraph{Lack of formal guarantees.}
We provide intuitive explanations for why reconstruction error contrast
indicates OOD status, but lack formal theoretical guarantees (e.g.,
PAC-style bounds). Claim~2 in
Section~\ref{sec:meth:diffusion:formal}
provides heuristic support but not a rigorous proof.

\paragraph{No adversarial evaluation.}
We do not evaluate robustness to adversarial examples crafted to fool the
OOD detector. Diffusion models have known vulnerabilities to adversarial
perturbations that preserve low-level statistics while crossing decision
boundaries.

%% ==================================================================
\section{Future Work}
\label{sec:disc:future}

\paragraph{Methodological extensions.}
Multi-class binary CDM evaluation (all 10 CIFAR-10 classes as ID in
turn); adaptive $\lambda$ selection via meta-learning or validation-set
AUROC; full-image or context-aware diffusion variants for inkjet,
leveraging latent diffusion~\cite{linmans2024diffusion} or layer-wise
reconstruction~\cite{yang2024diffusion}; crop conditioning that
incorporates spatial metadata.

\paragraph{Scaling and efficiency.}
Score distillation to a lightweight discriminative
model~\cite{wu2024diffusion}; progressive coarse-to-fine evaluation
($K{=}1$ triage, $K{=}50$ for borderline cases); extension to higher
resolutions ($256{\times}256{+}$) for medical and industrial settings.

\paragraph{Theoretical investigations.}
Formal bounds relating $\lambda$ to AUROC improvement; adversarial
robustness evaluation; application of reconstruction-error scoring to
text, audio, and time-series modalities.
\endgroup
